%%================================================
%% Chapter 1
%%================================================
\chapter{บทนำ}
\label{chapter1}

\section{ที่มาและความสำคัญ}
\label{sec:ch1_significance}

จากสถานการณ์คนพิการในประเทศไทย ปี พ.ศ. 2567 พบว่ามีประชากรผู้พิการกว่า 2.2 ล้านคน โดยเฉพาะผู้พิการทางการเคลื่อนไหว \cite{MSDHS_Disability_2024} ภาครัฐจึงได้มีนโยบายส่งเสริมคุณภาพชีวิตหลายด้าน ทั้งการสร้างอาชีพผ่านรายงานการจ้างงานคนพิการ \cite{DEP_Employment_2024} การสนับสนุนสิทธิและการใช้ชีวิตอิสระตามอนุสัญญาว่าด้วยสิทธิคนพิการ (Convention on the Rights of Persons with Disabilities: CRPD) \cite{UN_CRPD_2008} รวมถึงการประยุกต์ใช้เทคโนโลยีเพื่อการดำรงชีวิตอิสระสำหรับผู้สูงอายุและผู้พิการในแนวคิด Ambient Assisted Living (AAL) \cite{Rashidi2013AALSurvey} เพื่อลดอุปสรรคในการดำเนินชีวิต นอกจากนี้ยังมีการส่งเสริมการมีส่วนร่วมทางสังคมตามแผนพัฒนาการกีฬาแห่งชาติ ฉบับที่ 6 \cite{MOTS_SportsPlan_2022} และการสร้างสภาพแวดล้อมที่ปราศจากอุปสรรคตามแผนพัฒนาคุณภาพชีวิตคนพิการแห่งชาติ ฉบับที่ 5 \cite{DEP_DisabilityPlan_2023}

ผู้ใช้เก้าอี้รถเข็นและผู้สูงอายุในสถานดูแล เช่น บ้านพักคนชราวาสนะเวศม์ ซึ่งเป็นบริบทศึกษาของโครงงานนี้ ต้องพึ่งพาผู้ดูแลและเผชิญข้อจำกัดด้านการเคลื่อนที่และความปลอดภัย การติดตามตำแหน่ง การเคลื่อนไหว และสภาพร่างกายเบื้องต้นจึงเป็นข้อมูลสำคัญต่อการดูแลประจำวัน โดยเฉพาะเมื่อผู้พักอาศัยอยู่กระจายกันในหลายพื้นที่และผู้ดูแลไม่สามารถมองเห็นทุกจุดได้พร้อมกันตลอดเวลา

\begin{figure}[htbp]
    \centering
    \includegraphics[width=0.75\textwidth]{chapter1/ch1-fig04.jpg}
    \captionsetup{justification=centering}
    \caption{บริบทการดูแลผู้สูงอายุและผู้พิการที่ใช้เก้าอี้รถเข็นในสถานดูแลของประเทศไทย}
    \label{fig:ch1-fig01-thai-context}
\end{figure}

จากการลงพื้นที่สำรวจอาคารและพื้นที่บริการหลัก พบว่ามีผู้พักอาศัยกระจายอยู่ในหลายอาคาร ทำให้ผู้ดูแลไม่สามารถทราบตำแหน่งและสภาพของผู้ใช้เก้าอี้รถเข็นทุกคนได้ทันเวลา ซึ่งอาจนำไปสู่อุบัติเหตุ เช่น การหกล้ม หรือการไม่ได้รับความช่วยเหลือเมื่อเกิดเหตุผิดปกติ \cite{Sterling2016FallsInCareHomes}

\begin{figure}[htbp]
    \centering
    \includegraphics[width=0.75\textwidth]{chapter1/ch1-fig02.jpg}
    \captionsetup{justification=centering}
    \caption{อาคารและพื้นที่บริการหลักของบ้านพักคนชราวาสนะเวศม์ที่ใช้เป็นบริบทศึกษา}
    \label{fig:ch1-fig02-facility-context}
\end{figure}

นอกจากนี้ การทำงานของผู้ดูแลยังมีข้อจำกัดจากการบันทึกข้อมูลแบบแอนะล็อก เช่น กระดานหรือกระดาษ ซึ่งเสี่ยงต่อการสูญหายและไม่สามารถนำมาวิเคราะห์แนวโน้มสุขภาพได้อย่างต่อเนื่อง ปริมาณงานที่มากและการแจ้งเตือนที่ขาดบริบทประกอบยังอาจนำไปสู่ความล้าจากการแจ้งเตือน (alert fatigue) และทำให้การจัดลำดับความสำคัญของงานยากขึ้น \cite{Gates2018AlertFatigue}

\begin{figure}[htbp]
    \centering
    \includegraphics[width=0.75\textwidth]{chapter1/ch1-fig03.jpg}
    \captionsetup{justification=centering}
    \caption{การใช้กระดานบันทึกข้อมูลแบบแอนะล็อกที่พบในหน้างานจริง}
    \label{fig:ch1-fig03-analog-records}
\end{figure}

แม้ปัจจุบันจะมีระบบอัตโนมัติในบ้านและระบบระบุตำแหน่งภายในอาคารหลายรูปแบบ \cite{Home_Automation_Integration_2023,Indoor_GPS_Limitations,UWB_Accuracy_Spec} แต่มักไม่สอดคล้องกับพฤติกรรมของผู้ใช้เก้าอี้รถเข็น มีความซับซ้อน หรือมีต้นทุนสูง โครงงานนี้จึงเสนอการพัฒนาแพลตฟอร์ม WheelSense ซึ่งเป็นระบบสภาพแวดล้อมอัจฉริยะแบบบูรณาการที่เชื่อมโยงอุปกรณ์บนเก้าอี้รถเข็น อุปกรณ์สวมใส่ โหนดภายในอาคาร และเซิร์ฟเวอร์ขอบ เพื่อสนับสนุนการทำงานของผู้ดูแลและเพิ่มความปลอดภัยให้แก่ผู้ใช้งานอย่างมีประสิทธิภาพ โดยระบบยังคงอยู่ในระดับต้นแบบเพื่อสนับสนุนการเฝ้าระวังเบื้องต้น ไม่ใช่ระบบวินิจฉัยหรือเครื่องมือแพทย์ที่ได้รับการรับรอง

\section{วัตถุประสงค์}
\label{sec:ch1_objectives}

\begin{enumerate}
    \item เพื่อพัฒนาระบบต้นแบบสภาพแวดล้อมอัจฉริยะสำหรับผู้ใช้เก้าอี้รถเข็นที่สามารถรับข้อมูลการเคลื่อนไหว สัญญาณชีพ และบริบทของพื้นที่ภายในอาคารได้อย่างต่อเนื่อง
    \item เพื่อพัฒนาระบบระบุตำแหน่งภายในอาคารระดับโซนหรือห้อง โดยอาศัยสัญญาณ BLE และการประมวลผลบนเซิร์ฟเวอร์
    \item เพื่อพัฒนาแพลตฟอร์มซอฟต์แวร์สำหรับจัดเก็บ ประมวลผล และแสดงผลข้อมูลแก่ผู้ใช้งานตามบทบาท ผ่านเว็บ dashboard และแอปพลิเคชันบนอุปกรณ์เคลื่อนที่
    \item เพื่อบูรณาการผู้ช่วยปัญญาประดิษฐ์ภายใต้กรอบเสนอ-ยืนยัน-ดำเนินการ เพื่อช่วยสรุปข้อมูลและสนับสนุนการตัดสินใจของผู้ดูแลอย่างปลอดภัย
    \item เพื่อประเมินระบบต้นแบบในด้านความถูกต้องของการระบุตำแหน่ง ความเสถียรของการไหลของข้อมูล ประสิทธิภาพของส่วนบริการ และประสบการณ์ผู้ใช้ของผู้ดูแล
\end{enumerate}

\section{ขอบเขตงานวิจัย}
\label{sec:ch1_scope}

\begin{enumerate}
    \item พัฒนาชุดอุปกรณ์ภาคสนามสำหรับติดตั้งบนเก้าอี้รถเข็นและอุปกรณ์สวมใส่ที่ใช้เก็บข้อมูลการเคลื่อนไหวและสัญญาณชีพพื้นฐาน
    \item พัฒนาระบบระบุตำแหน่งภายในอาคารระดับโซนหรือห้อง โดยใช้ข้อมูล BLE RSSI และข้อมูลอ้างอิงจากพื้นที่ที่กำหนด
    \item พัฒนาเว็บ dashboard และแอปพลิเคชันบนอุปกรณ์เคลื่อนที่สำหรับการใช้งานตามบทบาทของผู้ดูแล ผู้ปฏิบัติงาน และผู้เกี่ยวข้อง
    \item บูรณาการผู้ช่วยปัญญาประดิษฐ์แบบภายในระบบ โดยจำกัดการกระทำที่มีผลเปลี่ยนสถานะให้อยู่ภายใต้ขั้นตอนยืนยันก่อนดำเนินการ
    \item ประเมินระบบต้นแบบภายใต้สภาพแวดล้อมทดลองและการใช้งานสาธิตในบริบทสถานดูแล โดยไม่มุ่งหมายเป็นการทดลองทางคลินิกหรือการรับรองเป็นเครื่องมือแพทย์
\end{enumerate}

\section{ประโยชน์ที่คาดว่าจะได้รับ}
\label{sec:ch1_benefits}

\textbf{ประโยชน์ต่อสังคม}
\begin{enumerate}
    \item ช่วยยกระดับคุณภาพการดูแลผู้สูงอายุและผู้ใช้เก้าอี้รถเข็นในสถานดูแล โดยเพิ่มความสามารถในการเฝ้าระวัง ความปลอดภัย การติดตามตำแหน่ง และการรับรู้เหตุผิดปกติได้รวดเร็วยิ่งขึ้น
    \item ช่วยลดภาระงานซ้ำซ้อนของผู้ดูแลและสนับสนุนการประสานงานข้ามเวรด้วยข้อมูลที่ค้นย้อนกลับได้ ซึ่งเอื้อต่อการจัดการทรัพยากรบุคคลอย่างมีประสิทธิภาพมากขึ้น
\end{enumerate}

\textbf{องค์ความรู้ที่ได้รับ}
\begin{enumerate}
    \item ได้องค์ความรู้ด้านการออกแบบระบบสภาพแวดล้อมอัจฉริยะที่บูรณาการอุปกรณ์ภาคสนาม บริการส่วนกลาง และส่วนติดต่อผู้ใช้เข้าด้วยกันในบริบทการดูแล
    \item ได้องค์ความรู้ด้านการประยุกต์ใช้ผู้ช่วยปัญญาประดิษฐ์อย่างปลอดภัยในงานเชิงปฏิบัติการที่เกี่ยวข้องกับข้อมูลผู้พักอาศัยและการตัดสินใจของผู้ดูแล
\end{enumerate}

\clearpage
\section{แผนการดำเนินงาน}
\label{sec:ch1_timeline}

การดำเนินงานวิจัยครอบคลุมระยะเวลา 12 เดือน โดยแบ่งเป็นช่วงการวิเคราะห์ปัญหา การศึกษาเอกสารที่เกี่ยวข้อง การออกแบบและสร้างระบบต้นแบบ การบูรณาการและทดสอบ ตลอดจนการวิเคราะห์ผลและจัดทำเล่มวิทยานิพนธ์ ดังรายละเอียดในตารางที่~\ref{tab:chapter1_timeline}

\begin{table}[H]
    \centering
    \caption{แผนการดำเนินงานของโครงงาน}
    \label{tab:chapter1_timeline}
    \small
    \renewcommand{\arraystretch}{1.35}
    \setlength{\tabcolsep}{2.2pt}
    \begin{tabular}{|L{5.8cm}|C{0.45cm}|C{0.45cm}|C{0.45cm}|C{0.45cm}|C{0.45cm}|C{0.45cm}|C{0.45cm}|C{0.45cm}|C{0.45cm}|C{0.45cm}|C{0.45cm}|C{0.45cm}|}
        \hline
        \multirow{2}{*}{\centering \textbf{กิจกรรม}} & \multicolumn{12}{c|}{\textbf{เดือนที่}} \\ \cline{2-13}
        & \textbf{1} & \textbf{2} & \textbf{3} & \textbf{4} & \textbf{5} & \textbf{6} & \textbf{7} & \textbf{8} & \textbf{9} & \textbf{10} & \textbf{11} & \textbf{12} \\ \hline
        1. วิเคราะห์ปัญหาและเก็บบริบทหน้างาน & \(\to\) & \(\to\) & & & & & & & & & & \\ \hline
        2. ศึกษาวรรณกรรมและทฤษฎีที่เกี่ยวข้อง & \(\to\) & \(\to\) & \(\to\) & & & & & & & & & \\ \hline
        3. ออกแบบสถาปัตยกรรมระบบและ data flow & & \(\to\) & \(\to\) & \(\to\) & & & & & & & & \\ \hline
        4. พัฒนา wheelchair sensor และ wearable gateway & & & \(\to\) & \(\to\) & \(\to\) & \(\to\) & & & & & & \\ \hline
        5. พัฒนาระบบระบุตำแหน่งภายในอาคาร & & & & \(\to\) & \(\to\) & \(\to\) & \(\to\) & & & & & \\ \hline
        6. พัฒนา backend, database และ MQTT service & & & & \(\to\) & \(\to\) & \(\to\) & \(\to\) & \(\to\) & & & & \\ \hline
        7. พัฒนา dashboard และ HMI ตามบทบาทผู้ใช้ & & & & & \(\to\) & \(\to\) & \(\to\) & \(\to\) & \(\to\) & & & \\ \hline
        8. บูรณาการ AI/MCP pipeline และ execution gate & & & & & & \(\to\) & \(\to\) & \(\to\) & \(\to\) & & & \\ \hline
        9. ทดสอบ sensor, localization, gateway และ server & & & & & & & & \(\to\) & \(\to\) & \(\to\) & & \\ \hline
        10. ประเมิน dashboard, AI และสรุปผลการทดลอง & & & & & & & & & \(\to\) & \(\to\) & \(\to\) & \\ \hline
        11. จัดทำและปรับปรุงปริญญานิพนธ์ & & & & & & & & & & \(\to\) & \(\to\) & \(\to\) \\ \hline
    \end{tabular}
    \\
    \vspace{0.5em}
    \footnotesize \textbf{หมายเหตุ:} \(\to\) หมายถึงช่วงเวลาที่ดำเนินกิจกรรม
\end{table}
